\textbf{COLUNI 2018 -- Questão 26}\par\medskip Leia a notícia, a seguir, que apresenta uma breve descrição de um acidente com graves impactos ambientais.

\par

Em julho de 2012, uma carreta que transportava aproximadamente 23 mil litros de ácido sulfúrico (\(\mathrm{H_2SO_4}\)) tombou na rodovia BR-381, no município de Nova Era, na região central de Minas Gerais. Cerca de 15 mil litros desse ácido atingiram a galeria pluvial e alcançaram o leito do rio Piracicaba. A população ribeirinha foi orientada a não utilizar a água do rio e, em seguida, observou-se a morte de espécies de peixes. CISCATO, C. A. M. et al. \textbf{Química}. v. 2. São Paulo: Moderna, 2016, p. 37. Adaptado.

\par

Para amenizar os danos ambientais em acidentes desse tipo, pode ser realizada a neutralização da quantidade de ácido presente no ambiente pela adição de soda cáustica (NaOH). Representa-se esse processo pela equação química não balanceada:

\par

\[\mathrm{x\ H_2SO_4(aq) + y\ NaOH(aq)} \longrightarrow \mathrm{z\ Na_2SO_4(aq) + w\ H_2O(l)}\]

\par

Com base nessas informações e nos conhecimentos de Ciências, assinale a afirmativa CORRETA.\begin{enumerate}
\item A substância que contaminou o rio, nesse acidente, é um exemplo de substância simples.
\item A neutralização de um ácido ocorre com a adição de um sal, formando uma base e água.
\item A neutralização do ácido com a adição da soda cáustica é um fenômeno físico.
\item Os coeficientes x, y, z e w que tornam a equação balanceada são, respectivamente, 1, 2, 1 e 2.\end{enumerate}
